% ------------------------------------------------------------------------
% Tablas del capítulo de RESULTADOS.
%
% Rellenadas con los valores medidos hasta la última corrida de
% preentrenamiento, ajuste fino y entrenamiento del decodificador.
%
% Lo que sigue con \dots NO se ha medido todavía. Cada hueco lleva encima un
% comentario con qué falta y de dónde sale. No se ha inventado ningún valor:
% en particular la tabla B3 queda vacía porque requiere cifras publicadas por
% terceros, que deben leerse de sus artículos y no reconstruirse de memoria.
%
% Paquetes que estas tablas necesitan:
%   \usepackage{booktabs, multirow, tabularx}
% ------------------------------------------------------------------------

% ---------------------------------------------------------------------
%  BLOQUE A · Configuración experimental
% ---------------------------------------------------------------------

% A1 · Conjuntos de evaluación
\begin{table}[htbp]
  \centering
  \caption{Conjuntos de regresión experimental: propiedad medida, unidades,
           tamaño y criterio de partición.}
  \label{tab:a1-conjuntos}
  \small
  \begin{tabular}{@{}llrrl@{}}
    \toprule
    Conjunto & Propiedad (unidades) & Moléculas & Test & Partición \\
    \midrule
    ESOL (Delaney)   & Solubilidad acuosa (log mol/L)        & 1\,128 & 113 & Aleatoria 80/10/10 \\
    FreeSolv (SAMPL) & Energía libre de hidratación (kcal/mol) &   642 &  65 & Aleatoria 80/10/10 \\
    Lipophilicity    & logD a pH 7,4 (adimensional)          & 4\,200 & 420 & Aleatoria 80/10/10 \\
    \bottomrule
  \end{tabular}

  \vspace{0.5em}
  \begin{minipage}{0.92\textwidth}
    \footnotesize
    La partición se rehace en cada semilla sobre el conjunto recombinado; los tamaños
    son constantes y la pertenencia varía. Los duplicados por identificador estructural
    no se fusionan, de modo que ESOL conserva las 1\,128 filas con que se publicaron las
    cifras de referencia.
  \end{minipage}
\end{table}


% A2 · Corpus y control de fuga
\begin{table}[htbp]
  \centering
  \caption{Control de fuga: moléculas de cada conjunto de evaluación y
           eliminadas del corpus de preentrenamiento por coincidencia de
           identificador estructural.}
  \label{tab:a2-fuga}
  \small
  \begin{tabular}{@{}lrrr@{}}
    \toprule
    Conjunto & Filas & Claves únicas & Eliminadas del corpus \\
    \midrule
    ESOL          & 1\,128 & 1\,117 &  37 \\
    FreeSolv      &    642 &    642 &   0 \\
    Lipophilicity & 4\,200 & 4\,200 & 169 \\
    \midrule
    \textbf{Total} & \textbf{5\,970} & \textbf{5\,564} & \textbf{196} \\
    \bottomrule
  \end{tabular}

  \vspace{0.5em}
  \begin{minipage}{0.92\textwidth}
    \footnotesize
    Solape de esqueletos reportado y conservado: 728 esqueletos compartidos,
    que cubren 304\,951 moléculas del corpus (15,7\,\%).
    El total de claves únicas no es la suma de las tres filas porque 395 moléculas
    figuran en más de un conjunto. La exclusión cubre los tres subconjuntos de cada
    dataset, no solo el de prueba, para que el corpus no dependa de la partición en uso.
  \end{minipage}
\end{table}


% A3 · Cobertura de distribución
\begin{table}[htbp]
  \centering
  \caption{Cobertura de distribución. Átomos pesados (percentiles 5, 50 y 95)
           y peso molecular mediano, sobre muestras de 20\,000 moléculas.}
  \label{tab:a3-cobertura}
  \small
  \begin{tabular}{@{}lrrrr@{}}
    \toprule
    & \multicolumn{3}{c}{Átomos pesados} & \\
    \cmidrule(lr){2-4}
    Conjunto & p5 & p50 & p95 & MW mediana \\
    \midrule
    MOSES         & 17 & 21 & 25 & 300 \\
    ZINC-250k     & 18 & 22 & 25 & 315 \\
    \midrule
    FreeSolv      &  4 &  8 & 18 & 120 \\
    ESOL          &  4 & 12 & 25 & 184 \\
    Lipophilicity & 15 & 27 & 38 & 388 \\
    \bottomrule
  \end{tabular}

  \vspace{0.5em}
  \begin{minipage}{0.92\textwidth}
    \footnotesize
    El corpus ocupa la banda 17--25 y falla por los dos extremos: FreeSolv queda
    íntegramente por debajo, la mediana de ESOL cae bajo su percentil 5 y más de la
    mitad de Lipophilicity supera su percentil 95. ZINC-250k ocupa esa misma banda, de
    modo que escalar el corpus con más ZINC añadiría volumen y ninguna cobertura nueva.
  \end{minipage}
\end{table}


% A4 · Equiparación de capacidad entre ramas
\begin{table}[htbp]
  \centering
  \caption{Equiparación de capacidad entre las dos ramas del codificador,
           con tres capas y ancho oculto de 256.}
  \label{tab:a4-capacidad}
  \small
  \begin{tabular}{@{}lrr@{}}
    \toprule
    Rama & Parámetros iniciales & Parámetros equiparados \\
    \midrule
    SMILES & 3\,945\,216 & 3\,945\,216 \\
    Grafo  &   450\,816 & 3\,219\,456 \\
    \midrule
    \textbf{Razón} & \textbf{8,75}$\times$ & \textbf{1,23}$\times$ \\
    \bottomrule
  \end{tabular}

  \vspace{0.5em}
  \begin{minipage}{0.92\textwidth}
    \footnotesize
    El desequilibrio inicial procede de que la capa transformadora fija su bloque
    interno en 2\,048 unidades sin escalarlo con el ancho oculto, mientras que el
    perceptrón de actualización del grafo operaba de ancho oculto a ancho oculto. La
    equiparación otorga a ambas ramas el mismo bloque de propagación hacia adelante; el
    1,23$\times$ restante corresponde al bloque de atención, cuya contrapartida en la
    rama de grafo es el paso de mensajes y no parámetros.
  \end{minipage}
\end{table}


% ---------------------------------------------------------------------
%  BLOQUE B · Mitad predictiva
% ---------------------------------------------------------------------

% B1 · Efecto del preentrenamiento  (tabla principal de la mitad predictiva)
\begin{table}[htbp]
  \centering
  \caption{Efecto del preentrenamiento sobre MOSES. Tres semillas por celda,
           contraste pareado sobre las mismas particiones. Las dos filas de
           cada conjunto comparten arquitectura y difieren solo en el origen
           de los pesos iniciales.}
  \label{tab:b1-preentrenamiento}
  \small
  \begin{tabular}{@{}llrrrrc@{}}
    \toprule
    Conjunto & Pesos iniciales & RMSE & MAE & $R^2$ & $\Delta$ pareada & Semillas \\
    \midrule
    \multirow{2}{*}{ESOL}
      & desde cero    & 0,741 $\pm$ 0,061 & 0,573 & 0,864 &                    &     \\
      & preentrenado  & \textbf{0,611} $\pm$ 0,040 & 0,459 & 0,907 & $-0{,}130 \pm 0{,}067$ & 3/3 \\
    \midrule
    \multirow{2}{*}{FreeSolv}
      & desde cero    & 1,420 $\pm$ 0,286 & 1,005 & 0,872 &                    &     \\
      & preentrenado  & \textbf{1,139} $\pm$ 0,136 & 0,707 & 0,918 & $-0{,}280 \pm 0{,}191$ & 3/3 \\
    \midrule
    \multirow{2}{*}{Lipophilicity}
      & desde cero    & 0,845 $\pm$ 0,063 & 0,625 & 0,526 &                    &     \\
      & preentrenado  & \textbf{0,671} $\pm$ 0,026 & 0,504 & 0,702 & $-0{,}174 \pm 0{,}054$ & 3/3 \\
    \bottomrule
  \end{tabular}

  \vspace{0.5em}
  \begin{minipage}{0.92\textwidth}
    \footnotesize
    Las nueve semillas favorecen a la fila preentrenada. Bajo la hipótesis nula de
    ausencia de efecto, nueve de nueve sobre particiones independientes tiene
    probabilidad $0{,}5^9 = 0{,}002$; por conjunto separado, tres de tres alcanza solo
    $0{,}125$, de modo que la afirmación se formula sobre las tres tareas conjuntamente.
    La dispersión de la columna $\Delta$ es la del contraste pareado, no la diferencia
    de las dispersiones de cada fila.
    \textbf{Medido con los duplicados aún fusionados} (ESOL con 1\,117 filas); la
    repetición sobre el conjunto de la tabla~\ref{tab:a1-conjuntos} está pendiente.
  \end{minipage}
\end{table}
% B3 · Contra la literatura publicada
%
% VACÍA DELIBERADAMENTE. Requiere cifras publicadas por terceros, que han de leerse de
% los artículos correspondientes. Reconstruirlas de memoria produciría valores plausibles
% y falsos, que es el peor resultado posible en una tabla de comparación.
%
% Qué anotar de cada trabajo que se cite, sin lo cual la fila es inutilizable:
%   - criterio de partición (aleatoria o por esqueleto) y proporciones
%   - número de semillas y si la partición se rehace en cada una
%   - si se fusionaron duplicados (cambia el tamaño de ESOL de 1 128 a 1 117)
%   - si el RMSE se reporta en unidades originales o estandarizadas
% Candidatos habituales: el artículo original de MoleculeNet, y trabajos de
% preentrenamiento molecular que reporten los tres conjuntos bajo partición aleatoria.
\begin{table}[htbp]
  \centering
  \caption{Situación frente a la literatura publicada. Cada fila indica el
           protocolo de partición bajo el que se obtuvo la cifra, sin el cual
           los valores no son comparables.}
  \label{tab:b3-literatura}
  \small
  \begin{tabularx}{\textwidth}{@{}Xlrrr@{}}
    \toprule
    Trabajo & Partición & ESOL & FreeSolv & Lipophilicity \\
    \midrule
    \dots\ \cite{}          & \dots & \dots & \dots & \dots \\
    \dots\ \cite{}          & \dots & \dots & \dots & \dots \\
    \dots\ \cite{}          & \dots & \dots & \dots & \dots \\
    \midrule
    \textbf{Este trabajo}   & Aleatoria, rehecha por semilla & \textbf{0,611} & \textbf{1,139} & \textbf{0,671} \\
    \bottomrule
  \end{tabularx}
\end{table}
% ---------------------------------------------------------------------
%  BLOQUE C · Mitad generativa
% ---------------------------------------------------------------------

% C1 · Evaluación con oráculo  (tabla principal de la tesis)
\begin{table}[htbp]
  \centering
  \caption{Evaluación generativa con oráculo. Cada brazo se evalúa con dos
           pesos de guía sobre las mismas 100 moléculas semilla del conjunto
           de prueba.}
  \label{tab:c1-oraculo}
  \setlength{\tabcolsep}{3.5pt}
  \footnotesize
  \begin{tabular}{@{}lrrrrrrrrr@{}}
    \toprule
    Brazo & $w$ & $\rho$ & Pend. & Valid. & Unic. & Nov. & Copia & Tanim. & Control \\
    \midrule
    \multirow{2}{*}{Fusión, desde cero}
      & 1 & $+0{,}811$ & $+0{,}900$ & 1,000 & 0,930 & 0,854 & 0,040 & 0,520 & $-0{,}263$ \\
      & 3 & $+0{,}856$ & $+1{,}145$ & 1,000 & 0,950 & 0,874 & 0,038 & 0,490 & $-0{,}181$ \\
    \addlinespace
    \multirow{2}{*}{Fusión, preentrenado}
      & 1 & $+0{,}778$ & $+0{,}893$ & 1,000 & 0,958 & 0,885 & 0,025 & 0,463 & $-0{,}222$ \\
      & 3 & $+0{,}806$ & $+1{,}158$ & 1,000 & 0,951 & 0,904 & 0,014 & 0,435 & $-0{,}222$ \\
    \addlinespace
    \multirow{2}{*}{Solo grafo}
      & 1 & $+0{,}786$ & $+0{,}866$ & 1,000 & 0,935 & 0,857 & 0,039 & 0,522 & $-0{,}145$ \\
      & 3 & $+0{,}851$ & $+1{,}114$ & 1,000 & 0,959 & 0,877 & 0,035 & 0,482 & $-0{,}153$ \\
    \addlinespace
    \multirow{2}{*}{Solo SMILES}
      & 1 & $+0{,}812$ & $+0{,}885$ & 1,000 & 0,933 & 0,845 & 0,044 & 0,525 & $-0{,}156$ \\
      & 3 & $+0{,}871$ & $+1{,}131$ & 1,000 & 0,950 & 0,876 & 0,034 & 0,481 & $-0{,}168$ \\
    \bottomrule
  \end{tabular}

  \vspace{0.4em}
  \begin{minipage}{0.95\textwidth}
    \footnotesize
    \textit{Fusión}: grafo y SMILES conjuntamente. \textit{Pend.}: pendiente de la
    recta que enfrenta el desplazamiento solicitado con el obtenido.
    \textit{Valid.}, \textit{Unic.}, \textit{Nov.}: validez, unicidad y novedad.
    \textit{Tanim.}: similitud de Tanimoto respecto de la molécula semilla.
    \textit{Control}: desplazamiento obtenido con la condición puesta al
    intervalo nulo; su desviación típica va de 0,638 a 0,783, frente a 0,3--0,5 por
    intervalo solicitado, de modo que condicionar desplaza la media y estrecha la
    dispersión. La validez es 1,000 por construcción de la representación empleada en
    la salida del decodificador, no por mérito del modelo.
  \end{minipage}
\end{table}


% C2 · Pedido contra obtenido, por intervalo
%
% INCOMPLETA: faltan los intervalos 0 a 9, que quedaron fuera del registro consultado.
% Están en results/oracle/pairs_graph_smiles_scratch_s2025_logp/summary.json, campo
% guidance["1.0"]["per_bin"], claves bin, requested_centre, obtained_mean, obtained_std,
% hit_rate, copy_rate, mean_tanimoto.
\begin{table}[htbp]
  \centering
  \caption{Calibración por intervalo del brazo fusionado sin preentrenar, con
           peso de guía unitario: desplazamiento solicitado frente al
           realmente obtenido.}
  \label{tab:c2-calibracion}
  \setlength{\tabcolsep}{5pt}
  \footnotesize
  \begin{tabular}{@{}rrrrrrr@{}}
    \toprule
    Intervalo & $\Delta$ solicitado & $\Delta$ obtenido & Desv. & Tasa de acierto & Copia & Tanimoto \\
    \midrule
    0  & \dots & \dots & \dots & \dots & \dots & \dots \\
    1  & \dots & \dots & \dots & \dots & \dots & \dots \\
    2  & \dots & \dots & \dots & \dots & \dots & \dots \\
    3  & \dots & \dots & \dots & \dots & \dots & \dots \\
    4  & \dots & \dots & \dots & \dots & \dots & \dots \\
    5  & \dots & \dots & \dots & \dots & \dots & \dots \\
    6  & \dots & \dots & \dots & \dots & \dots & \dots \\
    7  & \dots & \dots & \dots & \dots & \dots & \dots \\
    8  & \dots & \dots & \dots & \dots & \dots & \dots \\
    9  & \dots & \dots & \dots & \dots & \dots & \dots \\
    10 & $+0{,}007$ & $-0{,}048$ & 0,314 & 0,020 & 0,570 & 0,821 \\
    11 & $+0{,}069$ & $-0{,}003$ & 0,398 & 0,190 & 0,010 & 0,489 \\
    12 & $+0{,}166$ & $+0{,}099$ & 0,378 & 0,170 & 0,040 & 0,515 \\
    13 & $+0{,}255$ & $+0{,}294$ & 0,328 & 0,300 & 0,010 & 0,516 \\
    14 & $+0{,}347$ & $+0{,}434$ & 0,414 & 0,200 & 0,010 & 0,471 \\
    15 & $+0{,}452$ & $+0{,}446$ & 0,506 & 0,210 & 0,000 & 0,458 \\
    16 & $+0{,}580$ & $+0{,}662$ & 0,298 & 0,160 & 0,000 & 0,451 \\
    17 & $+0{,}738$ & $+0{,}684$ & 0,590 & 0,320 & 0,000 & 0,425 \\
    18 & $+0{,}985$ & $+0{,}899$ & 0,689 & 0,300 & 0,000 & 0,426 \\
    19 & $+1{,}296$ & $+1{,}296$ & 0,681 & 0,720 & 0,000 & 0,380 \\
    \bottomrule
  \end{tabular}

  \vspace{0.5em}
  \begin{minipage}{0.92\textwidth}
    \footnotesize
    El intervalo 10 corresponde a solicitar un desplazamiento nulo, y es el único con
    una tasa de copia alta (0,570) y un Tanimoto elevado (0,821): devolver la molécula
    de partida es una respuesta correcta a esa petición. Los intervalos primero y último
    son abiertos, de modo que su centro es una extrapolación y su error absoluto no es
    comparable con el de los centrales; en ellos debe leerse la tasa de acierto. Los
    intervalos centrales son estrechos por construcción, al proceder de veinte cuantiles
    de la distribución de desplazamientos, por lo que una tasa de acierto baja allí
    refleja la anchura del intervalo tanto como el comportamiento del modelo.
  \end{minipage}
\end{table}


% C3 · Ablación de los dos ejes
\begin{table}[htbp]
  \centering
  \caption{Ablación sobre los dos ejes: modalidades del codificador y origen
           de los pesos. Correlación entre desplazamiento solicitado y obtenido.}
  \label{tab:c3-ablacion}
  \small
  \begin{tabular}{@{}llrr@{}}
    \toprule
    & & \multicolumn{2}{c}{$\rho$} \\
    \cmidrule(lr){3-4}
    Eje & Brazo & $w = 1$ & $w = 3$ \\
    \midrule
    \multirow{3}{*}{Modalidad}
      & Grafo + SMILES & $+0{,}811$ & $+0{,}856$ \\
      & Solo grafo     & $+0{,}786$ & $+0{,}851$ \\
      & Solo SMILES    & $+0{,}812$ & $+0{,}871$ \\
    \midrule
    \multirow{2}{*}{Pesos iniciales}
      & Desde cero     & $+0{,}811$ & $+0{,}856$ \\
      & Preentrenado   & $+0{,}778$ & $+0{,}806$ \\
    \bottomrule
  \end{tabular}

  \vspace{0.5em}
  \begin{minipage}{0.92\textwidth}
    \footnotesize
    El brazo \emph{grafo + SMILES desde cero} es el pivote de ambos ejes y aparece dos
    veces. Los brazos difieren en 0,033 con $w = 1$ y en 0,065 con $w = 3$; con una
    única semilla de entrenamiento por brazo, ninguna de las dos separaciones es
    defendible como diferencia. La fusión multimodal no supera a ninguna de las
    modalidades por separado, y la inicialización desde el codificador preentrenado
    queda última en ambos pesos de guía.
  \end{minipage}
\end{table}


% C4 · Presupuesto y convergencia
\begin{table}[htbp]
  \centering
  \caption{Presupuesto de cómputo y convergencia de los cuatro brazos
           generativos, todos bajo el mismo número de pasos.}
  \label{tab:c4-presupuesto}
  \setlength{\tabcolsep}{5pt}
  \footnotesize
  \begin{tabular}{@{}lrrrrr@{}}
    \toprule
    Brazo & Parámetros & Pasos & Pérdida val. & Prec. tok. & Min. \\
    \midrule
    Grafo + SMILES, desde cero   & 14\,113\,319 & 60\,000 & 0,0953 & 0,9628 & 31 \\
    Grafo + SMILES, preentren.   & 14\,113\,319 & 60\,000 & 0,1249 & 0,9510 & 31 \\
    Solo grafo                   & 10\,135\,588 & 60\,000 & 0,1047 & 0,9588 & 16 \\
    Solo SMILES                  & 10\,893\,857 & 60\,000 & 0,0981 & 0,9616 & 26 \\
    \bottomrule
  \end{tabular}

  \vspace{0.5em}
  \begin{minipage}{0.92\textwidth}
    \footnotesize
    El presupuesto se fija en pasos y no en épocas para que un brazo cuyas épocas sean
    más baratas no entrene más tiempo y convierta la comparación en una sobre tiempo de
    reloj. La pérdida de validación mide reconstrucción de la secuencia de salida, no
    condicionamiento: dado que los pares minados comparten una similitud mínima
    apreciable, un modelo que copiase la molécula de partida puntuaría bien sin haber
    aprendido a condicionar. Se reporta para constatar convergencia, no para comparar
    brazos; esa comparación es la de la tabla~\ref{tab:c3-ablacion}.
  \end{minipage}
\end{table}
% D2 · Estadística de los pares minados
\begin{table}[htbp]
  \centering
  \caption{Estadística de los pares minados. Los percentiles delimitan el
           rango de desplazamiento que la evaluación puede solicitar.}
  \label{tab:d2-pares}
  \small
  \begin{tabular}{@{}lrrrr@{}}
    \toprule
    Propiedad & Desv. típica de $\Delta$ & p1 & p99 & $|\Delta| > 1\sigma$ \\
    \midrule
    logP & 0,685  & $-1{,}766$ & $+1{,}766$ & 27,7\,\% \\
    TPSA & 15,784 & $-43{,}09$ & $+43{,}09$ & 29,4\,\% \\
    QED  & 0,066  & $-0{,}198$ & $+0{,}198$ & 21,9\,\% \\
    MW   & 27,679 & $-69{,}09$ & $+69{,}09$ & 32,1\,\% \\
    \bottomrule
  \end{tabular}

  \vspace{0.5em}
  \begin{minipage}{0.92\textwidth}
    \footnotesize
    Pares totales: 12\,991\,314, de los cuales 649\,566 son de identidad.
    Tanimoto intra-par: rango admitido $[0{,}50,\ 0{,}95]$; la mediana está pendiente de
    extraer de \texttt{data/moses/pairs\_meta.json}.
    La media de $\Delta$ es nula y los percentiles simétricos, lo que confirma que los
    pares se emiten en ambas direcciones. La última columna se expresa en unidades de
    desviación típica porque un umbral absoluto carece de sentido comparativo entre
    propiedades cuyos rangos difieren en varios órdenes de magnitud.
    El percentil 99 de logP, $\pm 1{,}77$, define el rango de control que la evaluación
    puede solicitar; solicitar más sería extrapolación.
  \end{minipage}
\end{table}
